\documentclass[12pt]{article}
\usepackage{amsmath}
\usepackage{amssymb}
\usepackage{xcolor}
\usepackage[margin=1in]{geometry}
\usepackage{amssymb}
\usepackage[T1]{fontenc}

\title{Cognitive Fusion Mechanism: Integration of LLM and Memory}
\author{Witold Warcho\l{}}
\date{\today}

\begin{document}

\maketitle

\section{Overview}

The fusion mechanism projects two independent information streams
into a shared 64 dimensional space: the LLM embedding query and episodic
memory entries. The design prioritizes information preservation and prevents
any single stream from drowning out the other through deterministic,
non-learnable feature extraction followed by banded assembly and light
context-modulated mixing. Neuron graph reasoning is handled by the
Python-side \texttt{\_NeuronGraphReasoner} (GAT) and enters the fusion
head as a separate token sequence rather than through this C-side pipeline.

\section{Input Streams}

\subsection{LLM Query Stream}

The LLM provides a dense embedding vector:
\begin{equation}
\mathbf{q}_{\text{raw}} \in \mathbb{R}^{d_{\text{llm}}}
\end{equation}

where $d_{\text{llm}}$ is the model's embedding dimension (variable).
This stream represents the current model query or context.

\subsection{Memory Stream}

Episodic memory entries consist of a rich vector and an importance weight:
\begin{equation}
\mathbf{m}_i = (\mathbf{v}_i, \alpha_i, \tau_i)
\quad \text{where} \quad
\mathbf{v}_i \in \mathbb{R}^{2N + d_{\text{input}}}
\end{equation}

where $N = \texttt{MAX\_NEURONS}$ and $d_{\text{input}} = \texttt{INPUT\_SIZE}$.
Up to $M = 1200$ entries may be stored. The vector combines prior neuron
state (implicit history) and external input features.

\section{Projection Mechanism}

\subsection{Deterministic Random Projection}

Rather than learned weights, projections use a seeded pseudo-random
matrix fixed across runs. A projection coefficient is computed as:

\begin{equation}
P(r, c) = \frac{\text{bits}(h(r, c)) \& 0xffffff}{2^{23}} - 1.0
\end{equation}

where $h(r, c)$ applies splitmix64-style mixing:
\begin{align}
z &= (r \ll 32) \oplus (c \cdot 2654435761 + 1) \\
z &\leftarrow z + 0x9e3779b97f4a7c15 \\
z &\leftarrow (z \oplus (z \gg 30)) \times 0xbf58476d1ce4e5b9 \\
z &\leftarrow (z \oplus (z \gg 27)) \times 0x94d049bb133111eb \\
z &\leftarrow z \oplus (z \gg 31)
\end{align}

This ensures:
\begin{itemize}
\item Every input dimension feeds every output dimension (dense mixing)
\item No information is dropped via blocking or striding
\item The matrix is reproducible without storage
\item Variance is well-distributed across output dimensions
\end{itemize}

\subsection{Full Projection}

An $n$-dimensional source is projected into the 64 dimensional space:
\begin{equation}
\text{proj}(\mathbf{x})_i = \frac{1}{\sqrt{n}} \sum_{j=0}^{n-1}
x_j \cdot P(i, j)
\quad \forall i \in [0, 64)
\end{equation}

The $\frac{1}{\sqrt{n}}$ scaling follows Johnson-Lindenstrauss and
stabilizes variance regardless of source dimension.

\subsection{Banded Projection}

A source is projected into a contiguous band $[o, o + \ell)$ of the
output while leaving other positions untouched:
\begin{equation}
\text{proj\_band}(\mathbf{x}, o, \ell, \text{seed})_i =
\begin{cases}
\frac{1}{\sqrt{n}} \sum_{j=0}^{n-1} x_j \cdot P(i + \text{seed}, j)
&\text{if } i \in [o, o+\ell) \\
0 &\text{otherwise}
\end{cases}
\end{equation}

Each band uses a different seed (1009 for the query band, 3001 for the
memory band) so bands occupy distinct subspaces and do not share a
projection.

\section{Processing Pipeline}

\subsection{Stage 1: LLM Query with Text Integration}

The LLM embedding is projected and normalized:
\begin{equation}
\mathbf{q} = \text{LayerNorm}(\text{proj}(\mathbf{q}_{\text{raw}}))
\end{equation}

If an input sentence embedding is provided, it is projected separately,
normalized, and blended at 0.5 strength:
\begin{equation}
\mathbf{q} \leftarrow \text{LayerNorm}(\mathbf{q} +
0.5 \cdot \text{LayerNorm}(\text{proj}(\mathbf{t}_{\text{raw}})))
\end{equation}

This allows the input sentence to directly influence the query band
without being gated away by memory blending.

\subsection{Stage 2: Top-K Memory Attention}

Memory entries are scored by dot-product similarity to the query,
and only the top-$K$ entries ($K = \min(M, 32)$) are used. This prevents
softmax collapse to uniform weights when entry importances are similar.

\subsubsection{Blending Prior and Memory}

Each entry is weighted-blended with a default system prior:
\begin{equation}
\mathbf{w}_i = (1 - \beta) \mathbf{d} + \beta \mathbf{v}_i
\quad \text{scaled by} \quad
\tilde{\mathbf{w}}_i = \alpha_i \cdot \mathbf{w}_i
\end{equation}

where $\beta = \text{mem\_weight\_ratio} \in [0, 1]$ controls the blend
($\beta = 0$ uses pure default, $\beta = 1$ uses pure stored memory).

\subsubsection{Key Projection and Scoring}

Each blended vector is projected and scored:
\begin{equation}
\mathbf{k}_i = \text{LayerNorm}(\text{proj}(\tilde{\mathbf{w}}_i))
\quad,\quad
s_i = \frac{1}{\sqrt{64}} \langle \mathbf{q}, \mathbf{k}_i \rangle
\end{equation}

\subsubsection{Top-K Selection}

The top $K = \min(M, 32)$ entries are selected by score. Their scores
are normalized via softmax:
\begin{equation}
\pi_i = \text{softmax}(s_i) = \frac{\exp(s_i - \max_j s_j)}
{\sum_j \exp(s_j - \max_j s_j)}
\end{equation}

\subsubsection{Memory Vector Aggregation}

The memory contribution is the weighted sum:
\begin{equation}
\mathbf{m} = \text{LayerNorm}\left(\sum_{i \in \text{top-K}}
\pi_i \mathbf{k}_i \right)
\end{equation}

\subsection{Stage 3: Banded Assembly}

Each stream is re-projected into its own contiguous band:
\begin{align}
\text{fused}[0:32]  &= \text{proj\_band}(\mathbf{q},  0,  32, 1009) \\
\text{fused}[32:64] &= \text{proj\_band}(\mathbf{m}, 32, 32, 3001)
\end{align}

where $\text{BAND\_Q} = 32$, $\text{BAND\_M} = 32$, and
$\text{FUSION\_DIM} = \text{BAND\_Q} + \text{BAND\_M} = 64$.

\textbf{Key property:} Each stream owns a disjoint set of dimensions.
No gating mechanism controls how much of each stream appears; instead,
the downstream fusion transformer learns to weight the bands as needed.

\subsection{Stage 4: Context-Modulated Cross-Band Mixing}

A light mixing term is applied based on context:
\begin{equation}
\gamma = \text{sigmoid}(2 \cdot \text{context\_factor} - 1) \times 0.5
\end{equation}

For each dimension, a partner dimension is mixed in:
\begin{equation}
\text{cross}[i] = \text{fused}[i] + \gamma \cdot
\text{fused}[(i + 32) \bmod 64]
\end{equation}

This introduces low-amplitude interactions between bands without
allowing dominance of any stream.

\subsection{Stage 5: Output Normalization}

The final vector is layer-normalized:
\begin{equation}
\mathbf{out} = \text{LayerNorm}(\text{cross})
\end{equation}

\section{Layer Normalization}

Layer normalization standardizes a vector to zero mean and unit variance:
\begin{align}
\mu &= \frac{1}{n} \sum_{i=0}^{n-1} v_i \\
\sigma^2 &= \frac{1}{n} \sum_{i=0}^{n-1} (v_i - \mu)^2 \\
\text{LayerNorm}(\mathbf{v})_i &= \frac{v_i - \mu}{\sqrt{\sigma^2 + \epsilon}}
\end{align}

where $\epsilon = 10^{-6}$ prevents division by zero.

\section{Design Rationale}

\subsection{Why Deterministic Projection?}

The C implementation is a deterministic feature extractor, not a learner.
Since gradients are killed at the ctypes boundary, learnable mixing must
occur in the downstream Python fusion transformer. The C side's role is
to produce a rich, stable, non-degenerate feature vector.

\subsection{Why Banded Architecture?}

By carving the output into two disjoint bands, each stream is
guaranteed representation. The transformer downstream learns to weight
each band appropriately. Neuron graph information enters the fusion head
through the Python-side GAT token sequence, which preserves the full
graph topology rather than collapsing it into a fixed-size projection.

\subsection{Why Top-K Memory Attention?}

Attending over all $M$ memory entries with similar importances causes
softmax to collapse to uniform weights, reducing the memory vector to
noise averaging. Top-K selection ($K = 32$) ensures the attention
distribution is peaked and informative even at high memory capacities
($M = 1200$).

\subsection{Why Project Bands Separately?}

Each band uses a different seed in the pseudo-random projection.
This ensures the two streams occupy distinct subspaces and their
projected features do not interfere.

\section{Output Characteristics}

\begin{itemize}
\item \textbf{Dimension:} 64 ($= \text{BAND\_Q} + \text{BAND\_M} = 32 + 32$)
\item \textbf{Bands:} Query (dims 0--31), Memory (dims 32--63)
\item \textbf{Normalization:} Layer-normalized, mean 0, variance 1
\item \textbf{Mixing:} Light cross-band mixing at $\leq 0.5$ strength,
controlled by context
\item \textbf{Information:} Both C-side streams present in balanced
proportion; neuron information enters the fusion head separately via
GAT tokens
\end{itemize}

\end{document}
